%Je nach dem in welcher Sprache ihr euer Paper schreiben wollt, benutzt bitte entweder den Deutschen-Titel oder den Englischen (einfach aus- bzw. einkommentieren mittels '%')

%Deutsch
\section{Schlussfolgerungen und Ausblick}\label{sec:conclusion}

Der anhaltende Trend steigender Motorisierung erhöht den Druck auf städtische Infrastrukturen durch vermeidbaren Suchverkehr kontinuierlich und rückt technologische Lösungsansätze wie \ac{iot}-basierte \acp{sps} und serverlose Cloud-Architekturen in den Mittelpunkt. Für Parkhausbetreiber stellt sich angesichts des wachsenden Wettbewerbsdrucks zur effizienteren Auslastung ihrer Stellplätze dabei die zentrale Frage der wirtschaftlichen Sinnhaftigkeit: Können die notwendigen Investitionen in Sensorik und Cloud-Infrastruktur eines \ac{sps} durch den erwarteten wirtschaftlichen Nutzen gerechtfertigt werden? Die Diskussion verschiebt sich damit von der technischen Machbarkeit hin zur Suche nach belastbaren Kennzahlen für eine fundierte Investitionsentscheidung.

Um diese Lücke zu adressieren, stellt Drive \& Decide einen methodischen Ansatz vor, der auf einem prototypischen, \ac{iot}-basierten \ac{sps} mit einer serverlosen Architektur auf \ac{aws} aufbaut. Der Versuchsaufbau kombiniert ein 3D-gedrucktes Modell einer Parkgarage mit einer Cloud-Anbindung über \ac{mqtt}, \ac{aws} \ac{iot} Core, \ac{aws} Lambda und Amazon DynamoDB. Auf Basis dieses Prototyps wird ein Kosten-Nutzen-Modell aus Sicht eines Parkhausbetreibers entwickelt. Dieses stellt einmalige Investitionskosten (\ac{capex}) und laufende, nutzungsabhängige Cloud-Kosten (\ac{opex}) dem potenziellen Mehrerlös gegenüber.

Die wichtigsten Erkenntnisse dieses Papers sind, dass die Investition in ein \ac{sps} zwingend auf Kennzahlen wie \ac{tco}, \ac{roi} und dem Break-Even-Punkt basieren muss. Der durchgehend ereignisgetriebene Aufbau über serverlose Dienste ist dabei eine wesentliche Voraussetzung, um die laufenden Cloud-Kosten (\ac{opex}) belastbar zu erfassen und eine realistische Einschätzung der wirtschaftlichen Rentabilität zu ermöglichen.

Als Ausblick auf das Final Paper und die nächsten Arbeitsschritte wird der vorgestellte Versuchsaufbau physisch und virtuell umgesetzt. Das 3D-gedruckte Modell der Parkgarage wird mit Sensoren und Matchbox-Fahrzeugen ausgestattet, während parallel das \ac{aws}-Backend vollständig konfiguriert wird. Ziel ist es, anhand realer Messwerte und lastabhängiger Cloud-Kosten das entwickelte Kosten-Nutzen-Modell quantitativ zu evaluieren und Parkhausbetreibern eine empirisch fundierte Entscheidungsgrundlage zu liefern.
